\section{The Role of the Cabinet}

\subsection{Composition of the Cabinet}

The CIA represents the different centres of American power involved in understanding and responding to developments in Latin America in November 1975. Although the CIA constitutes the institutional core of the cabinet, its composition extends beyond intelligence officers to include figures from other institutions.

This diversity is essential to understanding the cabinet. Its members do not necessarily share the same priorities or interpretation of American interests. Intelligence officials may prioritize the collection and protection of information, while political figures may be more concerned with domestic legitimacy, diplomatic consequences, or the credibility of American foreign policy. Political figures, moreover, operate within an increasingly scrutinized political environment in which the limits of covert action have become a central question.

The presence of Orlando Letelier, a Chilean political figure and former diplomat living in exile, further broadened the cabinet’s perspectives. His position provides a contrasting view of the political developments taking place in Chile and the wider region, while introducing the interests and concerns of those opposing the military governments.

The CIA therefore represents the different institutions and individuals capable of influencing how the United States understands, manages, and responds to the situation in Latin America.

\subsection{Objectives}

The CIA’s central objective is to protect and advance U.S. strategic interests in Latin America while preserving the ability of the United States to respond to developments in the region. Containing the expansion of communist influence and maintaining the political stability of governments considered strategically important remain significant concerns, particularly in the context of the broader Cold War.

At the same time, the CIA must maintain access to reliable intelligence regarding governments, revolutionary organizations, and foreign involvement in the region. Understanding developments before they become strategically significant is essential to determining what response, if any, may be appropriate.

The CIA should also aim to preserve the credibility and freedom of action of the United States. The ongoing Church Committee investigations place American intelligence activities under unprecedented political scrutiny, making exposure a significant consideration in any decision involving covert activity. The cabinet must therefore balance strategic interests abroad with the legal, political, and diplomatic consequences of American involvement. It must be taken into consideration that no single approach necessarily protects American interests in every situation.

There are different ways in which the CIA can operate, making use of their resources and range of action. Its capabilities are extensive, but constrained by political considerations.

\subsubsection{Training and doctrine}

Officers from Chile's DINA, Argentina's SIDE, and other intelligence services pass through US-run training programs, including instruction in interrogation and counterinsurgency methods, at institutions such as the School of the Americas. This does not require the CIA to direct a single operation; it means the operational vocabulary and methods used are substantially shaped by prior American instruction.

\subsubsection{Communications and technical infrastructure}

The dictatorships’ entire operating premise, instant, cross-border sharing of target information, depends on secure communications. U.S.-supplied radio and communications equipment, provided originally for bilateral counterinsurgency cooperation, could plausibly be repurposed by member services to build a shared channel, without any single American official approving that specific use.

\subsubsection{Bilateral intelligence sharing}

The CIA maintains separate liaison relationships with each Southern Cone service individually. Information the Agency shares with, say, DINA for what it understands to be a narrow bilateral purpose could be pooled by DINA with intelligence partners, effectively laundering CIA-origin intelligence into the multilateral network without the Agency's knowledge or consent.

\subsubsection{Selective inaction}

Perhaps the most realistic and most dramatically useful lever for this committee: the CIA does not need to actively build Condor to be complicit in it. Simply declining to raise alarms, declining to brief Congress or the National Security Council on what station officers in Buenos Aires, Santiago, or Montevideo are picking up, or slow-walking a diplomatic démarche, are all forms of involvement through omission, and each is a decision this cabinet may face.

\subsubsection{Reporting versus tipping off targets}

The distinction between the CIA's job of watching operations and the separate, forbidden step of directing or resourcing is the exact line this committee will be asked to navigate, and it is also the line that in practice could prove very easy to blur.

\subsubsection{Exposure management}

Given the ongoing Church Committee hearings, a significant share of this cabinet's realistic work is not operational at all, it is about managing what Congress, the press, and foreign allies eventually learn, and when. Decisions about what to document, what to destroy, what to brief upward, and what to let remain deniable are all live options for the committee.

\subsection{Relationship to other Cabinets}

\subsubsection{Cabinet A - Regional Council}

The Latin American governments are important partners and sources of intelligence for the United States, particularly in the context of anti-communist cooperation. At the same time, their actions may create political and diplomatic risks for Washington. The CIA must assess the value of its relationships with individual governments and the possibility that their activities may exceed what the United States is prepared to support or acknowledge.

\subsubsection{Cabinet C - Revolutionaries}

Revolutionary movements represent a central focus of American intelligence collection in the region. Their activities, networks, and relationships with other actors may provide important information about the development of the conflict. Moreover, the existence of multiple revolutionary organizations with different objectives means that the threat cannot be understood as a single unified movement.

\subsubsection{Cabinet D - KGB}

The Soviet Union’s participation in the conflict poses the most direct threat against the U.S. Soviet intelligence represents a potential source of external influence over Latin American political developments. Information regarding Soviet activity may influence how Washington assesses revolutionary organizations and the governments facing them, as well as determining how they will obtain leverage over the Soviets and communism.
